% Capitulo3_Modelado.tex
\chapter{Modelado del Entorno y Modelo de Observación del Sensor}
\label{cap:modelado}

\section{Modelo de Observación del Sensor (Sensor Ideal de Nivel 1)}
Para el modelado de la capacidad de detección del vehículo aéreo no tripulado (UAV), se asume el modelo estándar de observación recomendado en la literatura de SAR denominado \textbf{Sensor Ideal de Huella Geométrica (Nivel 1)}.

Bajo esta formulación, la probabilidad de detección $P_d(x,y)$ de una víctima ubicada en las coordenadas espaciales $(x,y)$ dada la posición del dron $p_k = (x_k, y_k)$ a una altitud de vuelo $h$ se define mediante una función de masa determinista sin degradación probabilística:

\begin{equation}
P_d(x,y \mid p_k) = \begin{cases} 
1.0 & \text{si } \sqrt{(x - x_k)^2 + (y - y_k)^2} \le R_d \\ 
0.0 & \text{en otro caso} 
\end{cases}
\end{equation}

donde $R_d$ representa el radio de cobertura efectiva de la cámara u óptico a bordo, el cual viene determinado por la altitud $h$ y el ángulo del campo de visión (\textit{Field of View - FoV}) $\theta_{FoV}$:

\begin{equation}
R_d = h \cdot \tan\left(\frac{\theta_{FoV}}{2}\right)
\end{equation}

\subsection{Discretización del Dominio de Búsqueda}
El espacio continuo de la Casa de Campo de Madrid ($12.75\text{ km}^2$) se discretiza mediante una rejilla regular de celdas cuadradas de tamaño:

\begin{equation}
\Delta x = \Delta y = 10\text{ metros} \quad (\text{Superficie de celda } A_{celda} = 100\text{ m}^2)
\end{equation}

Si la proyección de la huella del sensor cubre el centroide de la celda de $10\text{ m}$, la celda se marca como observada y la probabilidad contenida en ella es absorbida por el mapa de creencias del agente.
